% !TEX program = lualatex
\documentclass[10pt,a4paper]{article}

\usepackage{fontspec}
\usepackage{polyglossia}
\setmainlanguage{russian}
\setotherlanguage{english}
\setmainfont{Times New Roman}
\setmonofont{Consolas}
\newfontfamily\cyrillicfonttt{Consolas}

\usepackage[a4paper,margin=20mm]{geometry}
\usepackage{microtype}
\usepackage{graphicx}
\usepackage{booktabs}
\usepackage{array}
\usepackage{caption}
\usepackage{enumitem}
\usepackage{amsmath}
\usepackage{hyperref}

\hypersetup{hidelinks}

\setlength{\parindent}{1.25cm}
\setlength{\parskip}{0pt}
\linespread{1.0}
\pagestyle{empty}
\captionsetup{font=small,labelsep=period}
\setlist{nosep}
\makeatletter
\renewcommand\section{\@startsection{section}{1}{\z@}%
  {-1.8ex plus -0.5ex minus -0.2ex}%
  {0.8ex plus 0.2ex}%
  {\normalfont\Large\bfseries}}
\makeatother

\begin{document}

\noindent УДК 004.415.53

\begin{center}
  {\bfseries А. А. Гладкий}\par
  Московский технический университет связи и информатики, г. Москва, Российская Федерация\par
  \href{mailto:gladkiy.a2004@gmail.com}{gladkiy.a2004@gmail.com}

  \vspace{0.8em}
  {\bfseries\MakeUppercase{Исследование методов управляемой инъекции отказов для тестирования отказоустойчивости распределённых приложений}}
\end{center}

\noindent\textbf{Аннотация.}
В работе сравниваются два способа выбора места для инъекции отказа. Первый выбирает конфигурацию из полного каталога, второй предварительно исключает точки, отсутствующие в контрольной трассе. Стенд состоит из трёх прикладных сервисов, PostgreSQL и Kafka. В него внесены четыре дефекта, связанные с повторным платежом, некорректными повторными попытками, потерей события и превышением времени ответа. Проведено по 30 серий для каждой стратегии при бюджете восемь экспериментов. Стратегия по трассе находила в среднем 4,0 дефекта против 2,6 у случайной стратегии; доля сработавших инъекций составила 100 и 60,% соответственно. Все четыре дефекта были повторно обнаружены в 40 контрольных запусках. Медианные накладные расходы слоя инъекции составили 2,53,%.

\noindent\textbf{Ключевые слова:}
инъекция отказов; распределённые приложения; отказоустойчивость; трассировка; граф вызовов; тестирование.

\begin{english}
\begin{center}
  {\bfseries A. A. Gladkiy}\par
  Moscow Technical University of Communications and Informatics, Moscow, Russian Federation\par
  \href{mailto:gladkiy.a2004@gmail.com}{gladkiy.a2004@gmail.com}

  \vspace{0.8em}
  {\bfseries A STUDY OF CONTROLLED FAULT INJECTION METHODS FOR TESTING THE RESILIENCE OF DISTRIBUTED APPLICATIONS}
\end{center}

\noindent\textbf{Abstract.}
The study compares two ways of selecting a fault injection point. The first draws a configuration from the complete catalogue, while the second removes points absent from a control trace. The testbed comprises three application services, PostgreSQL, and Kafka, with four seeded defects involving a duplicate payment, an incorrect retry policy, a lost event, and excessive response time. Thirty campaigns were run for each strategy with a budget of eight experiments. The trace-based strategy detected 4.0 defects on average versus 2.6 for random selection, while the shares of activated injections were 100% and 60%, respectively. All four defects were reproduced in 40 verification runs. The median overhead of the injection layer was 2.53%.

\noindent\textbf{Keywords:}
fault injection; distributed applications; resilience; tracing; call graph; testing.
\end{english}

\section{Введение}

Эксперимент по инъекции отказа задаётся точкой воздействия, типом сбоя и моментом его активации \cite{basiri2016}. Если выбранный компонент не участвует в проверяемом сценарии, запуск заканчивается так же, как контрольный, и почти не даёт сведений об отказоустойчивости системы. При случайном выборе такие запуски неизбежны: множество настроенных компонентов и связей шире пути, который фактически проходит запрос.

Распределённая трассировка даёт другой способ сформировать множество кандидатов. Из завершённых трасс можно восстановить граф вызовов и оставить в нём только узлы и связи, наблюдавшиеся при выполнении сценария. У этого подхода есть цена: сбор трасс создаёт нагрузку, а полученный граф не содержит редко выполняемых ветвей. Поэтому преимущество такой стратегии нельзя считать заранее установленным.

\section{Цель исследования}

Цель работы --- экспериментально сравнить случайный выбор точек инъекции с выбором по фактическому графу вызовов. Проверяемая гипотеза $H_1$ состоит в том, что при одинаковом бюджете запусков стратегия, использующая граф вызовов, обнаруживает внедрённые дефекты быстрее случайной. Нулевая гипотеза $H_0$ предполагает, что различия между стратегиями укладываются в случайный разброс результатов. За один эксперимент принимается выполнение фиксированного пользовательского сценария при одной активной инъекции.

\section{Материалы и методы}

Стенд развёртывался в Docker Compose и включал три сервиса на Python~3.12: сервис заказов, сервис платежей и обработчик уведомлений. Данные хранились в PostgreSQL~16, события передавались через Kafka~3.9.1. Трассы собирались OpenTelemetry Collector и сохранялись в Jaeger. Пользовательский сценарий состоял из создания заказа, регистрации платежа, перевода заказа в состояние \texttt{PAID}, публикации события \texttt{order.paid} и регистрации уведомления.

После каждого запуска проверялись пять инвариантов: отсутствие второго платежа для одного заказа; наличие события для оплаченного заказа; завершение запроса не позднее 2~с; соблюдение ограничения повторных попыток и запрета повтора постоянной ошибки; отсутствие повторного уведомления. Нарушение инварианта связывалось с одним из четырёх заранее внесённых дефектов (табл.~\ref{tab:defects}).

\begin{table}[ht]
  \centering
  \caption{Дефекты экспериментального стенда}
  \label{tab:defects}
  \begin{tabular}{c>{\raggedright\arraybackslash}p{0.67\textwidth}c}
    \toprule
    Код & Условие проявления & Инвариант \\
    \midrule
    B1 & потеря ответа после сохранения платежа и повтор запроса без ключа идемпотентности & I1 \\
    B2 & повтор запроса после постоянной ошибки HTTP~400 & I4 \\
    B3 & ошибка публикации в Kafka после фиксации состояния \texttt{PAID} & I2 \\
    B4 & недоступность платёжного сервиса в течение 3~с при дедлайне 2~с & I3 \\
    \bottomrule
  \end{tabular}
\end{table}

Точка инъекции задавалась стабильным идентификатором. Конфигурация также содержала тип отказа, задержку, код ответа и число допустимых активаций. Сервисы поддерживали четыре типа воздействия: задержку, ошибку, недоступность и потерю ответа после выполнения операции. При достижении точки в трассу добавлялись её идентификатор, тип отказа и номер активации. Перед очередным экспериментом таблицы PostgreSQL очищались, активная конфигурация сбрасывалась, а заказ получал новый идентификатор. Kafka и сервисы между экспериментами не перезапускались.

Сравниваются две стратегии:

\begin{enumerate}[label=\arabic*.,leftmargin=*]
  \item случайная стратегия выбирает конфигурации из полного каталога;
  \item стратегия по трассе выполняет контрольный сценарий, извлекает атрибуты \texttt{fault.point.id} из его трассы и исключает остальные точки.
\end{enumerate}

Полный каталог содержал 14 совместимых пар ``точка~--- тип отказа''. В контрольной трассе встретились пять точек, которым соответствовали восемь конфигураций каталога. В одной серии конфигурации выбирались без повторений. Для обеих стратегий использовались бюджет из восьми экспериментов и одно случайное зерно. Проведено 30 парных серий с зернами от 1 до 30; порядок стратегий чередовался между сериями.

Инъекция считалась полезной, если выбранная точка была достигнута и отказ активировался. Для каждой стратегии доля таких запусков вычислялась по формуле

\begin{equation}
  \text{Д}_{\text{полезн}} = \frac{N_{\text{полезн}}}{N_{\text{всего}}},
\end{equation}

где $N_{\text{полезн}}$ --- число сработавших инъекций, а $N_{\text{всего}}$ --- общее число запусков стратегии. Также фиксировались число найденных уникальных дефектов, номер эксперимента и время до первого дефекта. Парные результаты сравнивались односторонним критерием знаков; совпавшие пары из расчёта вероятности исключались.

Для проверки воспроизводимости каждая конфигурация B1--B4 запускалась десять раз. Накладные расходы измерялись в 100 парных запусках после пяти прогревочных. Сравнивались включённый слой инъекции без активного отказа и его обход. Распределённая трассировка работала в обоих режимах, поэтому измерение не включает её стоимость. Порядок режимов чередовался.

\section{Результаты и обсуждение}

\begin{table}[ht]
  \centering
  \caption{Результаты 30 серий при бюджете восемь экспериментов}
  \label{tab:results}
  \begin{tabular}{>{\raggedright\arraybackslash}p{0.44\textwidth}cc}
    \toprule
    Метрика & Случайная & По трассе \\
    \midrule
    Среднее число уникальных дефектов & 2,6 & 4,0 \\
    Медиана числа уникальных дефектов & 3 & 4 \\
    Медиана номера первого обнаружения & 2 & 1 \\
    Медиана времени до первого дефекта, мс & 1232 & 1082 \\
    Доля полезных инъекций, \% & 60 & 100 \\
    \bottomrule
  \end{tabular}
\end{table}

Стратегия по трассе обнаружила все четыре дефекта в каждой серии. Случайная стратегия нашла B1 в 66,7,\% серий, B2 в 60,0,\%, B3 в 63,3,\% и B4 в 70,0,\%. По числу уникальных дефектов стратегия по трассе дала лучший результат в 26 парных сериях, ещё в четырёх результаты совпали; односторонний критерий знаков дал $p=1{,}49\cdot10^{-8}$. По доле полезных инъекций улучшение наблюдалось во всех 30 парах ($p=9{,}31\cdot10^{-10}$).

Первый дефект стратегия по трассе находила раньше по числу экспериментов в 16 сериях, в 14 сериях номера совпали ($p=1{,}53\cdot10^{-5}$). Различие времени до первого дефекта меньше: стратегия по трассе была быстрее в 20 сериях, случайная --- в 10 ($p=0{,}049$). На эту метрику влияет тип первого выбранного отказа: B1 добавляет ожидание клиентского тайм-аута, а B4 --- трёхсекундную задержку.

Все 40 повторных запусков известных конфигураций снова обнаружили ожидаемый дефект. Полученная на стенде оценка воспроизводимости равна 100,\%. В замере накладных расходов медиана клиентского времени выросла с 26,43 до 27,10~мс, то есть на 2,53,\%. После удаления 5,\% крайних значений с каждой стороны среднее выросло на 3,45,\%. Медианная парная разность составила 0,97~мс. Значения 95-го процентиля равны 33,06 и 33,31~мс. Следовательно, на данном стенде проверка точек инъекции добавляет около одной миллисекунды, а редкие задержки среды сильнее влияют на обычное среднее.

Контрольная трасса содержала только точки основного сценария. Поэтому преимущество второй стратегии объясняется сокращением пространства поиска: восемь конфигураций вместо четырнадцати. Этот результат проверяет заявленную гипотезу, но не показывает, что трассировка всегда даёт полный каталог отказов. Если редкая ветвь не попала в контрольный запуск, расположенная в ней точка будет исключена.

Известный набор внедрённых дефектов упрощает проверку результатов, но ограничивает переносимость выводов. Эксперимент показывает эффективность поиска на выбранном стенде и сценарии; по нему нельзя оценить распространённость этих дефектов в промышленных системах. Таблицы базы очищались перед каждым запуском, но сервисы и Kafka не перезапускались. Влияние постепенного прогрева уменьшалось чередованием порядка стратегий между сериями.

\section{Заключение}

При одинаковом бюджете стратегия, ограниченная точками контрольной трассы, использовала все инъекции и во всех 30 сериях обнаружила четыре заложенных дефекта. Случайный выбор расходовал 40,\% запусков на точки, не участвовавшие в сценарии, и находил в среднем 2,6 дефекта. Результаты поддерживают гипотезу $H_1$ для исследованного стенда: предварительная фильтрация по фактическим взаимодействиям сокращает число экспериментов до обнаружения ошибки и увеличивает число найденных дефектов.

Практическая граница метода связана с полнотой контрольных трасс. Для сценариев с редкими ветвями потребуется несколько контрольных запусков или объединение трасс с другими источниками сведений о системе. Дальнейшая проверка должна включать больше пользовательских сценариев, динамическую топологию и стенды с большим числом сервисов.

\section*{Благодарности}

\begin{thebibliography}{9}

\bibitem{basiri2016}
Basiri A., Behnam N., de Rooij R., Hochstein L., Kosewski L., Reynolds J., Rosenthal C.
Chaos Engineering // IEEE Software. 2016. Vol.~33, no.~3. P.~35--41.
\href{https://doi.org/10.1109/MS.2016.60}{https://doi.org/10.1109/MS.2016.60}.

\end{thebibliography}

\section*{Сведения об авторе}

\noindent\textbf{Гладкий Андрей Антонович}, магистрант Московского технического университета связи и информатики (МТУСИ), Москва, Российская Федерация, \href{mailto:gladkiy.a2004@gmail.com}{gladkiy.a2004@gmail.com}.

\begin{english}
\section*{About the author}

\noindent\textbf{Andrey Antonovich Gladkiy}, master's student, Moscow Technical University of Communications and Informatics (MTUCI), Moscow, Russian Federation, \href{mailto:gladkiy.a2004@gmail.com}{gladkiy.a2004@gmail.com}.
\end{english}

\end{document}
